
Reinforcement Learning from Human Feedback (RLHF) has become the dominant paradigm for aligning language models to human preferences. Evaluation benchmarks measure preference agreement---how often humans prefer aligned model outputs over baseline alternatives---as the primary success metric. InstructGPT \citep{ouyang2022instructgpt} reports 85\% preference win rates versus GPT-3; Constitutional AI \citep{bai2022constitutional} demonstrates improved helpfulness and harmlessness scores through RLHF optimization.

However, preference agreement metrics measure only \textbf{unidirectional alignment} ($AI\rightarrow human)$: does the model produce outputs humans prefer? They do not measure \textbf{bidirectional alignment} ($human\rightarrow AI)$: do humans preserve critical evaluation capacity when exposed to aligned models? High preference agreement could indicate quality improvement on objective tasks (users genuinely prefer correct answers to arithmetic problems or factual questions) or preference homogenization on subjective tasks (users habituate to model style and converge on preferences for creative writing or opinion questions where diversity is legitimate). Existing metrics cannot distinguish these scenarios.

We propose \textbf{preference entropy}---Shannon entropy H = -$\Sigma$ p\_i log(p\_i) applied to human preference distributions---as an information-theoretic proxy for bidirectional alignment. Higher entropy indicates diverse judgments (preserved critical evaluation); lower entropy indicates convergence (potential homogenization). Unlike self-reported agency metrics requiring dedicated user surveys, preference entropy could be computed from existing RLHF benchmark data if the data structure supports it.

Our validation reveals a methodological limitation: standard RLHF datasets use pairwise comparison formats optimized for annotation cost, structurally incompatible with entropy measurement. Testing on Anthropic-HH \citep{bai2022constitutional}, we sampled n=100 prompts and computed Shannon entropy for aggregated preference distributions. Entropy computation succeeded technically (100\% success rate), but all prompts yielded identical H = ln(2) $\approx$ 0.693 nats (variance = 0)---not measurement noise, but a structural artifact.

\textbf{Root cause:} Each Anthropic-HH example compares different response candidates. Aggregation treats each pairwise comparison as ''1 vote chosen, 1 vote rejected,'' creating tautological 50/50 splits regardless of actual preference diversity. This pairwise-unique format enables efficient reward model training but prevents per-prompt diversity measurement. Multi-annotator formats (OpenAI Summarization; Stiennon et al., 2020) where multiple annotators rate the same response candidates would enable entropy measurement but incur higher annotation cost.

\subsubsection{Contributions}

\begin{enumerate}
\item \textbf{Bidirectional Alignment Framework (Proposed):} Introduce preference entropy as a lightweight proxy for collective critical evaluation capacity, measurable from preference data when structure supports it. First application of information-theoretic entropy to RLHF impact on human diversity assessment (proposal untested empirically).
\end{enumerate}

\begin{enumerate}
\item \textbf{Dataset Structure Documentation (Validated on n=1):} Document cost-measurement trade-off between pairwise-unique formats (Anthropic-HH: cost-efficient, entropy-incompatible, empirically validated) and multi-annotator-fixed formats (OpenAI Summarization: entropy-measurable, higher cost, format-based inference). Tested on Anthropic-HH (n=100 prompts); compatibility inference for WebGPT/InstructGPT based on documented formats.
\end{enumerate}

\begin{enumerate}
\item \textbf{Methodological Validation:} Demonstrate entropy computation is technically feasible (100\% success rate, correct mathematical range [0, ln(2)] for binary choices) but dataset format yields constant value (zero variance). Root cause analysis confirms structural incompatibility, not implementation error.
\end{enumerate}

\begin{enumerate}
\item \textbf{Alternative Proxies (Proposed, Untested):} Propose three alternatives when multi-annotator data unavailable: response diversity entropy (model outputs), intra-annotator variance (longitudinal data), entropy-regularized RLHF (algorithm modification). All proposals require future validation.
\end{enumerate}

\subsubsection{Key Finding}

Standard RLHF benchmarks (Anthropic-HH, WebGPT) cannot measure preference diversity even if desired. Pairwise comparison formats optimize annotation cost (1 labeler per comparison) but fundamentally prevent diversity measurement (tautological 50/50 aggregations). Future benchmarks must choose: train reward models efficiently (pairwise) OR evaluate bidirectional alignment (multi-annotator).

\textbf{Organization:} Section 2 reviews RLHF evaluation and bidirectional alignment gaps. Section 3 presents entropy methodology and dataset structure. Section 4 describes H-E1 validation experiment. Section 5 reports findings (constant entropy, root cause). Section 6 discusses implications and proposes alternatives. Section 7 concludes.

